\documentclass[conference]{IEEEtran}
\IEEEoverridecommandlockouts

% ── Packages ──────────────────────────────────────────────────────────────────
\usepackage{cite}
\usepackage{amsmath,amssymb,amsfonts}
\usepackage{algorithmic}
\usepackage{graphicx}
\usepackage{textcomp}
\usepackage{xcolor}
\usepackage{booktabs}
\usepackage{multirow}
\usepackage{array}
\usepackage{listings}
\usepackage{float}
\usepackage{hyperref}
\usepackage{url}
\usepackage{colortbl}
\usepackage{microtype}
\usepackage{balance}

% ── Colours ───────────────────────────────────────────────────────────────────
\definecolor{codebg}{RGB}{248,248,248}
\definecolor{coderule}{RGB}{200,200,200}
\definecolor{highlight}{RGB}{214,245,214}
\definecolor{accentblue}{RGB}{0,82,155}

% ── Hyperref ──────────────────────────────────────────────────────────────────
\hypersetup{
    colorlinks=true,
    linkcolor=accentblue,
    filecolor=accentblue,
    urlcolor=accentblue,
    citecolor=accentblue
}

\begin{document}

% ── Title ─────────────────────────────────────────────────────────────────────
\title{%
    \textbf{Apte: AI Principle Tracker Ethos}\\
    Automated Benchmarking and Comparative Analysis of \\
    Corporate AI Ethics Governance (2019--2025)%
}

\author{
    \IEEEauthorblockN{Souda Fathima Kurnool, Ali Akarsu, Jenil Jayeshbhai Parmar,}
    \IEEEauthorblockN{Smit Rajendrasinh Baghel, Dhruv Sharma}
    \IEEEauthorblockA{\textit{MSc in Artificial Intelligence}\\
    National College of Ireland, Dublin\\
    \{24251445, 24304051, 24281883, 25117149, 24234915\}@student.ncirl.ie}
}

\maketitle

% ─────────────────────────────────────────────────────────────────────────────
\begin{abstract}
As Artificial Intelligence (AI) systems become deeply integrated into societal infrastructure, the governance frameworks guiding their development—corporate AI ethics policies—have proliferated. However, interpreting these policies remains a challenge for researchers and the public due to their disperse, longitudinal, and often ambiguous nature. This paper presents \textbf{Apte (AI Principle Tracker Ethos)}, a full-stack automated system designed for the longitudinal and comparative analysis of AI ethics governance. Apte integrates two major data modalities: (1) a multi-year ethics keyword corpus for Gemini (2019–2025) consisting of over 6,900 records, and (2) structured policy datasets for leading AI corporations including Google, Microsoft, IBM, and OpenAI. Our methodology employs a hybrid approach combining keyword-frequency longitudinal tracking with a side-by-side comparison engine powered by baseline compliance scoring and LLM-assisted summarization. Results indicate significant semantic shifts in corporate priorities over the last six years, moving from broad "fairness" principles to specific "safety" and "governance" controls. The system provides a transparent, reproducible benchmarking tool that bridges the gap between high-level ethical principles and operational policy evidence.
\end{abstract}

\section{Introduction}
The rapid advancement of Large Language Models (LLMs) and generative AI has moved AI ethics from academic debate to a critical requirement for corporate governance. Major technology providers now publish extensive "AI Principles" and "Responsible AI" reports to demonstrate their commitment to ethical development. However, these documents are often criticized for being "ethics washing"—symbolic gestures that lack operational accountability \cite{munn2022uselessness}. The fundamental problem lies in the difficulty of tracking policy evolution over time and comparing commitments across different providers in a structured, quantitative manner.

The project \textbf{Apte (AI Principle Tracker Ethos)} aims to address this lack of transparency by providing a production-ready dashboard for benchmarking corporate AI ethics. The novelty of this work lies in its automated approach to policy analysis, repurposing large keyword datasets into longitudinal "ethics signals" and providing a side-by-side comparison engine that scores companies across six key ethical pillars: Fairness, Transparency, Accountability, Privacy, Safety, and Governance.

The research questions guiding this project are:
\begin{itemize}
    \item \textbf{RQ1:} How have corporate AI ethics priorities shifted semantically between 2019 and 2025?
    \item \textbf{RQ2:} Can automated keyword-based scoring effectively distinguish the maturity of ethics governance between major AI providers?
    \item \textbf{RQ3:} To what extent can a decoupled full-stack architecture support real-time, explainable AI policy benchmarking for non-expert stakeholders?
\end{itemize}

This report details the methodology, technical implementation, and findings of the Apte system, providing clear evidence of its significance in the field of AI governance and oversight.

\section{Related Work}
The landscape of AI ethics is characterized by a "proliferation of principles" but a "poverty of practice" \cite{mittelstadt2019ethics}. Jobin et al. provided a seminal global inventory of AI ethics guidelines, identifying a global convergence on five key principles: transparency, justice/fairness, non-maleficence, responsibility, and privacy \cite{jobin2019landscape}. However, their work was a static cross-sectional analysis that did not provide a tool for longitudinal tracking.

Critical evaluations by Hagendorff \cite{hagendorff2020ethics} and Munn \cite{munn2022uselessness} suggest that while the number of guidelines has grown, their impact on actual engineering practice remains questionable. Hagendorff notes that most guidelines lack enforcement mechanisms, leading to a gap in accountability. Our project, Apte, positions itself differently by focusing on the \textit{evidence} of policy evolution. By tracking the frequency and context of ethics-related terms over seven years (2019–2025), we provide a quantitative proxy for corporate focus that goes beyond mere principle statements.

Furthermore, the work of Raji et al. on algorithmic auditing highlights the need for internal and external accountability mechanisms \cite{raji2020closing}. Apte contributes to this by providing an external, automated benchmarking tool that can be used by auditors and researchers to identify when corporate language shifts significantly, potentially signaling changes in internal risk appetite or regulatory compliance.

Finally, the role of codes of ethics in software development has been studied by McNamara et al., who found that traditional codes often fail to influence developer behavior \cite{mcnamara2018acm}. By surfacing these principles in a high-visibility dashboard, Apte aims to increase the "social pressure" and transparency surrounding corporate commitments, making them harder to ignore or change without notice.

\section{Methodology}
\subsection{Datasets and Preprocessing}
Apte utilizes two primary data sources:
\begin{enumerate}
    \item \textbf{Gemini Ethics Dataset (2019–2025):} A collection of over 6,900 keyword records extracted from Gemini's internal and public ethics documentation. Each record includes a word and its association with ethics categories (e.g., "bias" mapped to Fairness).
    \item \textbf{Corporate Policy Corpus:} Structured CSV data containing policy points for Google, Microsoft, IBM, Amazon, Tesla, and OpenAI. Each entry includes the policy text, the year of adoption, and the relevant ethics category.
\end{enumerate}

Preprocessing involves normalising company names, cleaning whitespace, and mapping raw keywords to a unified taxonomy of six ethical pillars. We implemented a custom `EthicsDataService` in Python to handle the parsing of these heterogeneous sources, ensuring that both keyword-frequency data and explicit policy points can be compared on the same timeline.

\subsection{System Architecture}
The system follows a decoupled client-server architecture:
\begin{itemize}
    \item \textbf{Backend (FastAPI):} Handles data processing, baseline scoring, and session management. It exposes a RESTful API for fetching company lists, timelines, and comparison results.
    \item \textbf{Frontend (React/Vite):} A modern, responsive UI built with TypeScript. It uses Recharts for visualizing keyword trends and a modular component structure for the Comparison Dashboard and Analysis Panel.
    \item \textbf{Services Layer:} Includes a `BaselineService` for keyword-coverage scoring and a `SessionService` for tracking user analysis history.
\end{itemize}

\subsection{Comparison Methodology (Person 4 Contribution)}
The side-by-side comparison engine implements a "Keyword Coverage Scoring" (KCS) algorithm. For any two selected companies, the system:
1. Identifies the set of unique ethical keywords present in their policies.
2. Categorises these keywords into the six pillars.
3. Calculates a "Compliance Score" based on the variety and density of keywords per category.
4. Generates a comparative summary highlighting which company leads in specific areas (e.g., "Microsoft leads in Privacy, while Google leads in Transparency").

\subsection{Workflow and Pipeline}
Figure 1 (represented in text) shows the full pipeline:
\texttt{Data Ingestion (CSV) -> Normalisation -> Category Mapping -> Baseline Scoring -> API Exposure -> UI Rendering (Timeline/Dashboard)}.

\section{Results and Evaluation}
\subsection{Longitudinal Findings: The Shift in Priorities}
Analysis of the Gemini dataset reveals a clear semantic evolution. In 2019, the most frequent keywords were related to "Fairness" and "Human Condition" (e.g., bias, inclusive). By 2024–2025, there is a marked increase in terms related to "Safety," "Risk Mitigation," and "Governance" (e.g., guardrails, red-teaming, compliance). This suggests a shift from aspirational principles to operational controls as AI systems reached production scale.

\subsection{Comparative Analysis Results}
Using the Comparison Dashboard, we benchmarked several leading providers. Table I summarizes the compliance scores derived from the system.

\begin{table}[h]
\caption{Baseline Ethics Compliance Scores (Scale 0-10)}
\begin{center}
\begin{tabular}{|l|c|c|c|c|c|}
\hline
\textbf{Company} & \textbf{Fairness} & \textbf{Transparency} & \textbf{Safety} & \textbf{Privacy} & \textbf{Overall} \\
\hline
Google & 8.5 & 9.0 & 7.5 & 8.0 & 8.25 \\
\hline
Microsoft & 8.0 & 8.5 & 8.5 & 9.0 & 8.50 \\
\hline
OpenAI & 7.0 & 7.5 & 9.0 & 7.0 & 7.63 \\
\hline
IBM & 7.5 & 8.0 & 7.0 & 8.5 & 7.75 \\
\hline
\end{tabular}
\end{center}
\end{table}

The results show that Microsoft and Google maintain high scores across the board, with Microsoft leading in Privacy governance. OpenAI, while newer, shows a specialized focus on Safety and Risk Mitigation, reflecting its "Frontier AI" positioning.

\subsection{System Performance and Scalability}
The FastAPI backend demonstrates sub-100ms response times for timeline generation, even with the 6,900-record dataset, due to efficient in-memory caching in the `EthicsDataService`. The decoupled architecture allows for easy extension; new company data can be added simply by dropping a CSV file into the `data/` directory without changing the codebase.

\section{Conclusions and Future Work}
Apte successfully demonstrates that automated systems can provide meaningful benchmarks for AI ethics governance. By integrating longitudinal keyword tracking with side-by-side comparisons, the project provides a comprehensive view of how corporate commitments evolve and how they compare to competitors.

\textbf{Significance and Impact:} This work provides a practical tool for "Social Oversight" of AI. It empowers researchers to move past the "symbolic compliance" phase and hold corporations accountable to their stated principles through quantitative evidence.

\textbf{Limitations and Future Work:} Currently, the system relies on static CSV datasets. Future versions will incorporate real-time web scraping of corporate policy pages and sentiment analysis of public reports to provide a more dynamic view. Additionally, we plan to integrate "Inter-rater Reliability" metrics for the user review framework to ensure the human-in-the-loop component is robust.

\begin{thebibliography}{99}
\bibitem{jobin2019landscape} A. Jobin, M. Ienca, and E. Vayena, ``The global landscape of AI ethics guidelines,'' \textit{Nature Machine Intelligence}, vol. 1, no. 9, pp. 389--399, 2019.
\bibitem{mittelstadt2019ethics} B. Mittelstadt, ``AI ethics – too principled to practice?,'' \textit{Nature Machine Intelligence}, vol. 1, no. 11, pp. 501--507, 2019.
\bibitem{floridi2018ai4people} L. Floridi et al., ``AI4People—An Ethical Framework for a Good AI Society: Opportunities, Risks, Principles, and Recommendations,'' \textit{Minds and Machines}, vol. 28, no. 4, pp. 689--707, 2018.
\bibitem{whittaker2018ainow} M. Whittaker et al., ``AI Now Report 2018,'' \textit{AI Now Institute}, 2018.
\bibitem{ieee2019ead} IEEE, ``Ethically Aligned Design: A Vision for Prioritizing Human Well-being with Autonomous and Intelligent Systems,'' \textit{IEEE}, 2019.
\bibitem{hagendorff2020ethics} T. Hagendorff, ``The Ethics of AI Ethics: An Evaluation of Guidelines,'' \textit{Minds and Machines}, vol. 30, no. 1, pp. 99--120, 2020.
\bibitem{boddington2017towards} P. Boddington, \textit{Towards a Code of Ethics for Artificial Intelligence}, Springer, 2017.
\bibitem{munn2022uselessness} L. Munn, ``The uselessness of AI ethics,'' \textit{AI and Ethics}, vol. 3, pp. 1--12, 2022.
\bibitem{raji2020closing} I. D. Raji et al., ``Closing the AI accountability gap: Defining challenges for internal algorithmic auditing,'' \textit{Proc. FAT* '20}, pp. 33--44, 2020.
\bibitem{schiff2020whats} D. Schiff et al., ``What’s Next for AI Ethics, Policy, and Governance? A Survey of Expert Expectations,'' \textit{arXiv preprint arXiv:2005.12652}, 2020.
\bibitem{mcnamara2018acm} A. McNamara, J. Smith, and E. Murphy-Hill, ``Does ACM's Code of Ethics Change Ethical Decision Making in Software Development?,'' \textit{Proc. ESEC/FSE '18}, pp. 729--733, 2018.
\end{thebibliography}

\end{document}
